\chapter{Lyme Disease Antibody Test}

\section*{What Is This Test?}
This test detects antibodies to the bacterium \textit{Borrelia burgdorferi}, the causative agent of Lyme disease. It utilizes a standardized two-tier testing protocol recommended by public health guidelines to ensure high diagnostic specificity \cite{cdc2023lyme}.

\section*{Alternative Names}
Lyme serology, Borrelia burgdorferi antibody, Lyme ELISA, Lyme Western Blot, two-tier Lyme test

\section*{Principle}
Diagnosis follows a two-tier algorithm \cite{lantos2021clinical}. The first tier is an enzyme-linked immunosorbent assay (ELISA) or immunofluorescence assay (IFA) measuring total IgM/IgG antibodies. If the first tier is positive or equivocal, it is followed by a second-tier Western Blot (immunoblot). The Western Blot detects antibodies against specific molecular weight proteins of the spirochete (IgM requires 2 of 3 bands: 23, 39, 41 kDa; IgG requires 5 of 10 bands: 18, 23, 28, 30, 39, 41, 45, 58, 66, 93 kDa).

\section*{History}
Lyme disease was first recognized in 1975 in Lyme, Connecticut \cite{steere2016lyme}. Dr. Willy Burgdorfer isolated the causative spirochete in 1981. The two-tier serological testing guidelines were formally standardized at the Dearborn Conference in 1994 to prevent false-positive diagnoses caused by cross-reactive antibodies.

\section*{Why This Test Is Ordered}
Your doctor has ordered this test if you exhibit symptoms of Lyme disease, such as an erythema migrans (bulls-eye) rash, joint pain, facial palsy, meningitis, or chronic fatigue, especially if you have a history of tick exposure in endemic regions.

\section*{Primary vs Secondary Test}
This is a primary screening and confirmatory diagnostic test for Lyme disease.

\section*{How This Test Is Performed}
A healthcare professional will draw a sample of blood from a vein in your arm.

\section*{What Preparation Is Needed}
No fasting or special preparation is required.

\section*{Contraindications and Precautions}
\begin{itemize}
  \item Standard precautions for venipuncture apply. The test should not be ordered as a general screening tool in patients without clinical symptoms or geographic exposure, as this increases the likelihood of false-positive results.
\end{itemize}
\section*{Risks}
There is minimal risk associated with a blood draw. You may experience minor bruising or discomfort at the puncture site.

\section*{What Happens After the Test}
Results are typically available in 3 to 7 days. If the two-tier test is positive and you have active symptoms, your doctor will prescribe a course of antibiotics (typically doxycycline or amoxicillin).

\section*{Understanding Results}

\subsection*{Below Normal}
A negative result (no antibodies detected) indicates that you likely do not have Lyme disease. However, early in the infection (the first 2 to 4 weeks), antibodies may not yet have reached detectable levels (the window period).

\subsection*{Above Normal}
A positive two-tier result indicates the presence of antibodies to \textit{Borrelia burgdorferi}. This confirms exposure and, in the presence of clinical symptoms, supports a diagnosis of active Lyme disease. It does not indicate active infection in asymptomatic individuals with past treated Lyme disease.

\section*{Normal Results}
Negative for both ELISA and Western Blot.

\section*{Follow-up Tests}
If symptoms persist despite a negative early test, repeat serology is recommended in 2 to 4 weeks. In patients with neurological symptoms, a lumbar puncture for cerebrospinal fluid (CSF) antibody index or PCR may be indicated.

\section*{References}
\bibliographystyle{unsrtnat}
\bibliography{references}

\clearpage
\section*{Regulatory Status and Authorization Date}
This chapter describes a test category, not a specific commercial product. FDA, CDSCO, EU IVDR, and MHRA approval or registration dates therefore cannot be assigned without the assay manufacturer, product name, instrument, and intended use. For laboratory-developed tests, an individual agency approval date may not exist. See the Preface for the interpretation of regulatory status.

\clearpage
